The confirmation experiment supports a bounded positive result: a model fitted before a severe crisis can improve response cost against static planning and a fixed capacity rule in the routing-feedback setting. A predictive selector often discards the relevant return dependency. Choosing actions across heterogeneous mechanisms also improves descriptive normalized cost against both fixed-action heuristics at the default cost, but the method is neither uniformly superior nor a guarantee of crisis prevention. Its losses without feedback, sensitivity to action cost, and reliance on supplied intervention effects remain material limits.
